\chapter{Introduction}

Reliable radio coverage prediction is a central requirement for future 6G
systems in which UAVs may act as aerial base stations, relays, or temporary
coverage nodes. A classical ray-tracing simulator can produce detailed channel
knowledge maps, but this is too expensive to repeat every time the transmitter
height, the city morphology, or the deployment assumptions change. In the 6G
literature, a channel knowledge map is a site-specific, location-tagged
repository of channel information that can support environment-aware planning
and reduce online channel-acquisition overhead
\cite{zengx2021ckm,ckmtutorial2024}. This thesis studies whether a learned
surrogate can predict dense channel maps directly from urban geometry while
preserving enough physical structure to generalize to cities not seen during
training.

The project is formulated as dense image-to-image regression on
$513\times513$ pixel city maps. The input is primarily a building-height
topology map, a LoS/NLoS mask, and the UAV antenna height. The outputs studied
through the project are path loss, delay spread, and angular spread. The final
line of work combines interpretable non-deep-learning priors with deep residual
models, so the neural network learns corrections around a strong physical
anchor instead of rediscovering the whole propagation law from scratch.

\section{Work goals}

The initial technical targets were defined in physical units and were kept as
the main success criteria throughout the project:

\begin{itemize}
  \item path-loss RMSE below \SI{5}{dB};
  \item delay-spread RMSE below \SI{50}{ns};
  \item angular-spread RMSE below \SI{20}{\degree};
  \item separate reporting for global, LoS, and NLoS regions.
\end{itemize}

The most important methodological objective was not only to obtain a low
average RMSE, but to do so under a strict city-holdout protocol. In this
setting, complete cities are reserved for validation and testing. This is a
harder and more deployment-relevant condition than random pixel or random tile
splits, because the model must extrapolate to urban layouts it has never seen.

\section{Requirements and specifications}

The project uses the CKM HDF5 dataset
\texttt{CKM\_Dataset\_270326.h5}. The canonical dataset contains 66 real-city
scenarios and 16180 samples. Each sample is a $513\times513$ pixel raster with
the following fields:

\begin{table}[ht]
\centering
\caption{Main CKM dataset fields used in the thesis.}
\label{tab:dataset_fields_intro}
\begin{tabularx}{\textwidth}{lllX}
\toprule
Field & Shape & Type & Role \\
\midrule
\texttt{topology\_map} & $513^2$ & float32 & Building height in metres. Pixels with buildings are not valid receiver locations. \\
\texttt{los\_mask} & $513^2$ & uint8 & Binary LoS/NLoS visibility mask. \\
\texttt{path\_loss} & $513^2$ & uint8 & Ray-traced path loss in dB, quantized at approximately \SI{1}{dB}. \\
\texttt{delay\_spread} & $513^2$ & uint16 & Delay-spread target, stored as integers and evaluated in ns after decoding. \\
\texttt{angular\_spread} & $513^2$ & uint8 & Angular-spread target, stored as integers and evaluated in degrees after decoding. \\
\texttt{uav\_height} & scalar & float32 & UAV transmitter height, used for conditioning and physical priors. \\
\bottomrule
\end{tabularx}
\end{table}

All losses and metrics are computed only on ground pixels. Building pixels are
kept in the input topology because they explain propagation, but they are
excluded from error accumulation because the receiver is assumed to be on the
ground. The fixed receiver height used in the analytic priors is
\SI{1.5}{m}. The transmitter is the UAV antenna, whose height varies
continuously across the dataset.

\section{Methods and procedures}

The work started from a U-Net and pix2pix-style cGAN baseline
\cite{isola2017pix2pix}. This was a natural first approach because the task is
dense map prediction, but the adversarial component did not solve the physical
errors that mattered most. Later attempts changed decoder upsampling, added
multi-scale supervision, separated metrics and regimes, introduced physical
prior channels, and finally moved to topology-partitioned experts.

The largest methodological change came after the models repeatedly failed on
NLoS tails. Histogram analysis showed that the path-loss distribution in NLoS
regions is not well described by a single Gaussian-like target. This motivated
Try 76: a distribution-first architecture with a sample-level GMM head followed
by a map head that assigns pixels to distribution components. Try 77 applied
the same idea to delay and angular spreads, where spike-plus-tail behaviour is
even more visible.

The final stage changed the role of deep learning again. Try 78 built a strong
interpretable path-loss prior from FSPL, coherent two-ray LoS propagation, and
calibrated NLoS corrections. Try 79 built log-domain ridge-calibrated priors
for delay and angular spread. Try 80 then combined these frozen priors with a
shared multi-task residual network.

\section{Work plan}
\label{sec:workplan}

The original work plan assumed a mostly linear path: literature review, dataset
preparation, baseline U-Net/cGAN training, model improvement, evaluation, and
writing. The actual project became more experimental because each stage exposed
a different limitation. Figure~\ref{fig:gantt} therefore shows the plan as the
major technical phases rather than as all 80 individual attempts.

\begin{figure}[ht]
  \centering
  \makebox[\textwidth][c]{%
\begin{ganttchart}[
  x unit=0.4cm,
  y unit title=0.6cm,
  y unit chart=0.6cm,
  vgrid,
  hgrid,
  title height=1,
  bar/.style={draw,fill=cyan!55},
  bar incomplete/.append style={fill=yellow!55},
  progress label text={},
  title label font=\small,
  group label font=\small\bfseries,
  bar label font=\small,
  milestone label font=\small\itshape,
  group/.style={draw,fill=gray!25},
  bar height=0.7,
  milestone/.style={fill=black}
]{1}{19}

  \gantttitle{Project Timeline (Weeks)}{19} \\
  \gantttitle{Jan}{2}
  \gantttitle{Feb}{4}
  \gantttitle{Mar}{4}
  \gantttitle{Apr}{5}
  \gantttitle{May}{4} \\

  \ganttgroup{WP1: Introduction and Research}{1}{5} \\
  \ganttbar[progress=100]{Proposal, workplan, SOA review}{1}{3} \\
  \ganttbar[progress=100]{Literature on 6G UAV formulas}{1}{5} \\

  \ganttgroup{WP2: Data Augmentation \& Heuristics}{8}{16} \\
  \ganttbar[progress=100]{Spread-target branch and first prior tests}{8}{10} \\
  \ganttbar[progress=100]{Ground-mask evaluation protocol}{10}{10} \\
  \ganttbar[progress=100]{Physical-prior residual baselines}{10}{12} \\
  \ganttbar[progress=100]{Frozen calibrated path-loss and spread priors}{14}{16} \\

  \ganttgroup{WP3: Deep Learning Pipeline}{4}{17} \\
  \ganttbar[progress=100]{Initial U-Net and cGAN baseline}{4}{6} \\
  \ganttbar[progress=100]{Height conditioning and visibility-aware splits}{6}{6} \\
  \ganttbar[progress=100]{Stabilized CNN inputs and distance features}{6}{7} \\
  \ganttbar[progress=100]{Artifact reduction and multiscale supervision}{7}{8} \\
  \ganttbar[progress=100]{PMHHNet and topology-expert consolidation}{11}{13} \\
  \ganttbar[progress=100]{Regularization and multi-expert diagnostics}{12}{14} \\
  \ganttbar[progress=100]{Distribution-first GMM modelling}{13}{14} \\
  \ganttbar[progress=100]{Joint prior-anchored residual GMM model}{16}{17} \\

  \ganttgroup{WP4: Testing}{14}{17} \\
  \ganttbar[progress=100]{Model evaluation against ground truth}{14}{17} \\

  \ganttgroup{WP5: Project Report}{6}{19} \\
  \ganttbar[progress=100]{Critical review and documentation}{6}{11} \\
  \ganttbar[progress=100]{Final thesis writing and formatting}{14}{19} \\
  \ganttmilestone{Final Review \& Presentation}{19}
\end{ganttchart}%
}

  \caption[Project's Gantt diagram]{Gantt diagram of the project, grouped by
  the main architectural phases rather than by individual attempts.}
  \label{fig:gantt}
\end{figure}

The main deviation from the initial plan was the discovery that a standard
image-to-image CNN could look plausible in LoS-dominated regions while still
failing badly in NLoS regions. This forced several changes: stricter ground-only
masking, regime-level diagnostics, topology experts, physical priors, and
finally the distribution-first and prior-anchored designs. The result is a
project whose final method is more hybrid and more interpretable than the
initial proposal anticipated.
